% abstract.tex — 摘要（素材与全部数字见 PAPER_HANDOFF.md, 2026-08-15）

生鲜商超的蔬菜商品保鲜期短、品相随销售时间下降，商超需在凌晨进货、以成本加成方式定价，
并受限于有限的销售空间。本文基于某商超 2020 年 7 月至 2023 年 6 月共 87.8 万条销售流水、
批发价格与损耗率数据，研究蔬菜类商品的自动定价与补货决策问题。

针对问题一，从总量、月度、周内与日内时段四个维度刻画六大品类销量的分布规律，发现花叶类
销量占三年总量的 \textbf{42.2\%}，周六周日与早晨时段的销量明显偏高，单品销量呈长尾分布；
进而以散点图检验结合斯皮尔曼秩相关分析品类间与单品间的关联关系，筛选出相关系数绝对值
大于 0.6 且显著性水平高于 0.01 的品类对 2 对、单品对 20 对，并结合上市季节与烹饪搭配
解释其成因。

针对问题二，先以三次多项式刻画各品类销售总量与成本加成定价的关系（拟合优度 $R^2$ 最高
0.58），再以历史 6--7 月同星期均值预测未来一周销量与批发价，构造含损耗品打折回收项的
利润函数，建立以补货量为决策变量的线性规划模型（方案 A）与价格离散化 0-1 混合整数规划
模型（方案 B），求得未来一周补货与定价策略，方案 A 七日总利润 \textbf{6735.03 元}，
且逐日利润高于 2021、2022 年同期水平。定价灵敏度分析表明，定价每变动 5\%，利润约同向
变动 16\%。

针对问题三，以利润额、销售量、销售次数、打折次数与损耗率为指标，用熵权法确定权重并通过
TOPSIS 对 49 种可售单品排序，取前 33 种作为候选；建立含各品类需求满足率不低于 50\%
约束的单品 0-1 规划模型，得到 7 月 1 日 30 个单品的补货与定价方案，单日利润
\textbf{1052.95 元}，各品类需求满足率达 \textbf{76.2\% 至 114.8\%}，兼顾利润目标与
市场需求的均衡。

针对问题四，从现有模型的简化假设出发，指出库存废弃、保鲜时间、天气客流、竞争定价与
供货约束五类数据缺口，对每类数据说明其采集方式、给出引入模型后的目标函数或约束修正
形式，并分析修正对补货与定价决策的改进效果。

全文模型经 \textbf{26 项一致性测试}验证：维度、手算用例、约束满足性、10 次重复运行
稳定性（标准差小于 $10^{-3}$）与交叉验证均通过；对空间容量与定价参数的灵敏度分析表明
模型结论稳健，支撑材料含全部代码与结果表。
